\documentclass{ximera}
\title{Logarithmic and Exponential Functions CW1}

\newcommand{\pskip}{\vskip 0.1 in}

\begin{document}
\begin{abstract}
Catenaries the number $e$.
\end{abstract}
\maketitle


\section{Compound Interest}

\begin{question} \label{QKf33r3rs}
\begin{enumerate}
\item Explain the meaning of the number $e$ in the context of this class and compound interest.

\item Write an expression that gives a good approximation to $e$ \emph{without} using $e$.

\end{enumerate}
\end{question}


\begin{question} \label{Qlerr33}
The balance in Account $A$ is $\$4,000$ at the start of $2026$. The account earns interest at the nominal rate of $8\%$/yr, compounded quarterly.

The balance in Account $B$ is $\$3,000$ at the start of $2026$. The account earns interest at the nominal rate of $12\%$/yr, compounded continuously.

Use the methods of \emph{this class} to determine exactly when the two accounts have the same balance. Then use a calculator to approximate this time.

Start by defining two functions that give the balances in the accounts. Define your inputs and outputs in complete sentences, with units.
\end{question}


\begin{question} \label{Qoero3r9030}
An initial investment of $\$5,000$ earns interest at the nominal rate of $8\%$/yr, compounded continuously.

\begin{enumerate}
\item How much interest does the investment earn during the tenth year?

\item Find a three-year period during which the investment earns $\$4,000$ in interest. Start by defining an unknown in a complete sentence.
\end{enumerate}
\end{question}


\begin{question} \label{Q6319L}
The balance in an account grows exponentially.

The account earns $\$1,000$ interest during the 10th year and $\$2000$ in interest during the 20th year. How much interest does it earn in 
\begin{enumerate}
\item the fifth year?

\item the third year?

\item in year $t$?
\end{enumerate} \end{question}


\section{Logistic Growth}

\begin{question} \label{Q33332}
Let
\[
   y =  f(x) = \frac{24}{x} .
\]

\begin{enumerate}
\item Describe the action of $f$ on its input.

\item Describe the action of $f^{-1}$ on its input.

\item Use your description in part (b) to find an expression for the function $x=f^{-1}(y)$.
\end{enumerate}

\end{question}

\begin{question} \label{Qdfl3r3r}
The function
\[
     P =f(t) = \frac{20}{1+3 \cdot 2^{-\frac{1}{10}t}} \, , \, 0\leq t \leq 75
\]
expresses the number of people (in thousands) infected with a virus in terms of the number of days since the start of the year. A graph of the function is shown below.

\begin{onlineOnly}
    \begin{center}
\desmos{1haivflyxu}{450}{600}  
\end{center}
\end{onlineOnly}

\href{https://www.desmos.com/calculator/1haivflyxu}{141: Logistic Model}

\begin{enumerate}
\item Is $f$ one-to-one? Explain.

\item Describe the action of $f$ on its input.

\item Describe the actionof $f^{-1}$ on its input.

\item Find an expression for the function $f^{-1}$. Use the appropriate variable names for the input and output.

\item When are 15,000 people infected with the virus? First use the graph to estimate the time. Then use your expression for $f^{-1}$ and a calculator to estimate the time.

\item Verify algebraically that $f^{-1}$ undoes $f$.
\end{enumerate}
\end{question}


\section{Hyperbolic Trig Functions}
\begin{question} \label{Qkdfsm13}
The graph of the hyperbolic cosine function
\[
    f(x) \cosh x= \frac{1}{2}\left( e^x + e^{-x} \right).
\]
is shown below.

\begin{onlineOnly}
    \begin{center}
\desmos{nubdwzgsa9}{450}{600}  
\end{center}
\end{onlineOnly}

\href{https://www.desmos.com/calculator/nubdwzgsa9}{152: Hyp Cosine}

\begin{enumerate}
\item Drag the slider $u$ in Line 2 of the worksheet above. Explain how to get the graph of the function $y=\cosh x$ from the graphs of the functions $y=e^x$ and $y=e^{-x}$.

\begin{freeResponse}
\end{freeResponse}

\item Is $f$ one-to-one? If not, what would be one way to restrict its domain to get a new function that is one-to-one and has the same range as $f$?

\item Find an expression for the inverse of the function
\[
    c(x) = \frac{1}{2}\left( e^x + e^{-x} \right).
\]

Start by making the substitution $u=e^x$ to get a quadratic equation in $u$. Then complete the square.



\item Activate the folder \emph{Hyperbolic Sine Function} in Line 15 above to see the graph of the hyperbolic sine function
\[
  g(x) =  \sinh(x) = \frac{1}{2}\left( e^x - e^{-x} \right).
\]

\item Is the function $g$ one-to-one?

\item Find an expression for the inverse of the function
\[
    g(x) = \frac{1}{2}\left( e^x + e^{-x} \right).
\]

Start by making the substitution $u=e^x$ to get a quadratic equation in $u$. Then complete the square.


\end{enumerate}

\end{question}


\section{Hanging Chains}

Hold the ends of a chain and let the chain sag under its own weight. The resulting curve looks like a parabola, and so it was thought. But it turns out that this is not correct. 


\begin{onlineOnly}
    \begin{center}
\desmos{tj3dz2cnf0}{450}{600}  
\end{center}
\end{onlineOnly}

\href{https://www.desmos.com/calculator/tj3dz2cnf0}{152: Hanging Chain 1}


\begin{enumerate} 

\item To see that the dashed (orange) chain in the worksheet above is \emph{not a parabola} drag the slider $c$ in Line 2.

\item It turns out that a chain with uniform density hanging in a uniform graviatational field assumes the shape of the hyperbolic cosine function $y=\cosh x$ (we'll see why later in the course). All that's needed to describe a particular chain is to scale the size of the graph by some factor. Test this by dragging the slider $a$ in Line 6 to make the graph of the function
\[
     f(x) = -a + a\cosh(x/a)
\]
match the highlighted (orange) chain below.
\end{enumerate}


\section{Some Questions}

\begin{question} \label{QLfeenvd}
\begin{enumerate}
\item Find an expression for the inverse of the function
\[
     f(x) = \cosh x = \frac{1}{2}\left( e^x + e^{-x} \right) \, , \, x\geq 0.
\]
Start by making the substutition $u=e^x$. Solve the resulting quadratic equation by completing the square.

\item Find an expression for the inverse of the function
\[
    g(x) = \sinh x = \frac{1}{2}\left( e^x + e^{-x} \right) \, , \, x\in \mathbb{R}.
\]
Start by making the substutition $u=e^x$. Solve the resulting quadratic equation by completing the square.

\item Simplify the composition
\[
       f(g^{-1}(x))
\]
as much as possible.
\end{enumerate}
\end{question}

\begin{question} \label{Qmdfmzmcp023}
The function
\[
    h = f(x) = h_0 + a\cosh(x/a) \, , \, x\geq 0,
\]
expresses the height (in feet) of a point on hanging chain in terms of its distance (in feet) from the chain's axis of symmetry.

The function
\[
      s = g(x) = a\sinh(x/a) \, , \, x\geq 0,
\]
gives the length of the chain (in feet) measured from its low point to a point on the chain $x$ feet from the chain's axis of symmetry.

\begin{enumerate}

\item What are the units of $h_0$? How do you know?

\item What are the units of $a$? How do you know?

\item Find a function
\[
     h = w(s) \, , \, s\geq 0,
\]
that expresses the height (measured in feet) of a point $P$ on the chain in terms of the length of the chain from its low point to $P$.

\end{enumerate}

\begin{onlineOnly}
    \begin{center}
\desmos{5oppqdy2l1}{450}{600}  
\end{center}
\end{onlineOnly}

\href{https://www.desmos.com/calculator/5oppqdy2l1}{141: Catenary}


\end{question}


\end{document}